\DFGHeading{objectives-work}

\DFGHeading{duration}
Illustrativer Zeitraum: 24 Monate.

\DFGHeading{objectives}
Das Beispiel zeigt Querverweise, Tabellen, Abbildungen und Literaturangaben.
\label{sec:objectives}

\DFGHeading{work-programme}
\DFGApplicantWork{pi1}
Die Ziele stehen in Abschnitt~\ref{sec:objectives}. Beide fiktiven Antragstellenden haben klar zugeordnete Aufgaben.
\DFGWorkPackage{WP1}{Methodenvergleich}
Alex Example entwickelt den Vergleich; Sam Example prüft die Reproduzierbarkeit.
\begin{equation}\label{eq:error} E = \frac{1}{n}\sum_{i=1}^{n}(y_i-\hat y_i)^2 .\end{equation}
\begin{figure}[htbp]\centering
\setlength{\unitlength}{1mm}\begin{picture}(130,22)
\put(0,5){\framebox(35,12){Material}}
\put(37,11){\vector(1,0){9}}\put(48,5){\framebox(35,12){Analyse}}
\put(85,11){\vector(1,0){9}}\put(96,5){\framebox(34,12){Prüfung}}
\end{picture}\caption{Fiktiver Arbeitsablauf zur Demonstration.}\label{fig:workflow}\end{figure}
\begin{table}[htbp]\centering\begin{tabularx}{\linewidth}{lXcc}\toprule
Aufgabe & Verantwortung & Monate 1--12 & Monate 13--24\\\midrule
WP1 & Alex Example & \cellcolor{black!12} X & \\
WP2 & Sam Example & & \cellcolor{black!12} X \\\bottomrule\end{tabularx}
\caption{Illustrativer Zeitplan.}\label{tab:schedule}\end{table}

\DFGHeading{data}
Hier würde die projektspezifische Begründung stehen. Im Beispiel werden keine Forschungsdaten erzeugt.

\DFGHeading{diversity}
Hier würde die projektspezifische Begründung stehen. Im Beispiel werden keine Forschungsdaten erzeugt.
